\section{Introduction}
Generative AI tools powered by Large Language Models (LLMs), such as GitHub Copilot and ChatGPT, are becoming increasingly integrated into software development workflows. These systems support capabilities including code completion, natural language-driven code generation, explanation, and debugging assistance \cite{vaithilingam2022expectation, barke2022grounded}. Recent studies report widespread adoption and measurable improvements in developer productivity and task completion speed \cite{peng2023copilot, bird2023copilot}. As a result, Generative AI is reshaping how developers interact with code and development environments.

Recent research has also begun examining how Generative AI influences developer behavior during programming tasks. Developers increasingly rely on prompt-based workflows, iterative refinement, and AI-generated code integration during implementation and debugging \cite{vaithilingam2022expectation, kazemitabaar2023copilot}. These interactions affect how solutions are explored, evaluated, and verified, while also introducing differences in trust and usage patterns across developers and tasks \cite{barke2022grounded, mao2026empirical}.

To better understand these changes, this paper presents a literature survey of ten studies on developer interaction with Generative AI tools. We first define a taxonomy of AI-assisted programming interactions, including suggestion-based, conversational, and generative workflows. Using this taxonomy, we analyze how these interactions influence stages of software development such as problem understanding, implementation, and debugging. By synthesizing findings across empirical studies, this work provides a structured view of how Generative AI is shaping developer behavior and programming workflows.

\subsection{Motivation}
The vision of an AI assistant collaborating with programmers has existed for decades across both programming language and machine learning research communities \cite{ferdowsifard2021LooPy,miltner2019edit,raychev2014code,xu2020externalknowledge}. Recent advances in Large Language Models (LLMs) have brought this vision closer to reality, with systems such as Codex and AlphaCode demonstrating strong code generation performance on programming tasks \cite{li2023competition,vaswani2023attention, chen2021evaluatingllm}. These capabilities have accelerated the adoption of industrial AI programming assistants such as GitHub Copilot.

Generative AI has emerged as a transformative paradigm capable of producing human-like text, code, and other artifacts with increasing fidelity. Modern LLMs demonstrate strong language understanding and generation abilities, enabling their adoption across a wide range of cognitive and collaborative tasks \cite{brown2020gpt3, bommasani2021foundation}. As a result, AI systems are increasingly transitioning from analytical tools to interactive assistants that support human problem solving.

In software development, tools such as GitHub Copilot and ChatGPT assist with programming tasks including code generation, explanation, and debugging. Studies report improvements in productivity, reduced task completion time, and support for learning unfamiliar concepts through AI-assisted workflows \cite{peng2023copilot, vaithilingam2022expectation, kazemitabaar2023copilot}. These systems also introduce new interaction patterns such as prompt-based programming and iterative refinement, reshaping how developers approach coding tasks.

At an industry level, Generative AI is influencing software engineering practices, productivity expectations, and developer workflows. Organizations increasingly adopt AI-assisted tools to improve efficiency and reduce development effort \cite{bird2023copilot, nocode2023industry}. However, there is still limited understanding of how these tools affect developer problem-solving strategies and cognitive processes, highlighting the need for structured analysis of AI-assisted programming behavior.

\subsection{Research Questions}
To better understand how Generative AI is reshaping software development practices, this paper investigates the influence of AI-assisted programming on developer behavior, problem-solving strategies, and software engineering workflows. Specifically, we structure our analysis around the following research questions:
\begin{enumerate}[label=\textbf{RQ \arabic*:},leftmargin=*]
    \item How do Generative AI tools reshape developers' problem solving strategies during software development?
    \item How does AI-assisted development influence developer behavior across different stages of the SDLC?
    \item What productivity benefits and tradeoffs are associated with AI-assisted software development?
\end{enumerate}